\section{Edwards Aquifer: the J-17 triggers under an affine recharge--pumping map}\label{sec:edwards}

\subsection{Data, map and declarations}\label{sec:edw-data}

The state is the annual mean elevation $H_t$ of the J-17 index well (TWDB well 6837203, ft above mean sea level), the control is total well pumping $P_t$ and the driver is total recharge $R_t$ (both in $10^3$ acre-ft; \cite{Umphres2025,EAA2023,TWDB}), on a locked 90-year panel 1934--2023 (Figure~\ref{fig:data}b). The map is the one-pool water balance in affine form,
\begin{equation}\label{eq:edw-map}
H_{t+1}-H_t=\alpha+\beta R_{t+1}+\gamma P_{t+1}+\delta H_t+e_{t+1},\qquad a=1+\delta,
\end{equation}
fitted by ordinary least squares on the 56 transitions with destination years 1935--1990: $\alpha=163.49$, $\beta=0.01983$, $\gamma=-0.02844$, $\delta=-0.2539$, so $a=0.7461$ and the map is a contraction. The sign pattern is the physical one (recharge raises, pumping lowers, head decays toward an equilibrium). The fit residuals have standard deviation $5.60$ ft, lag-one autocorrelation $-0.31$ and largest absolute value $15.41$ ft (1971), which is the declared defect. The 33 out-of-sample transitions 1991--2023 have standard deviation $8.40$ ft and largest absolute value $21.81$ ft (2004): the declaration is \emph{exceeded} out of sample, and Section~\ref{sec:edw-horizon} carries both defects. The floor classes are recharge floors from the fit window---the 1956 minimum $43.7$, the 5th percentile $166.5$ and the 10th percentile $179.1$ ($10^3$ acre-ft)---and the reference pumping is the fit-window mean $\Pbar=282.16$. Two thresholds are declared: 618 ft (the physical trigger of the Habitat Conservation Plan) and 660 ft (the Stage-I institutional trigger), with the ceiling at 710 ft. The rule family comprises BAU ($P=\Pbar$), flat caps at 90, 80, 70, 60, 50 and 0\% of $\Pbar$, the Stage-I rule ($-20\%$ below 660 ft) and the four-stage Critical Period Management cascade (cuts of 20, 30, 35 and 40\% below 660, 650, 640 and 630 ft).

\subsection{Attractors, kernels and the securability of the two triggers}\label{sec:edw-kernels}

By Corollary~\ref{cor:edw} every constant-pumping question is decided by the worst-floor attractor $\Hstar=(\alpha+\beta\underline R+\gamma P)/(1-a)$, and by Proposition~\ref{prop:stage} a staged rule is decided by the attractor of the stage that contains it. Table~\ref{tab:edw} lists them. Under the 1956 recharge floor BAU settles at $615.7$ ft, below 618; every cut of 10\% or more settles above it (flat 90\%: $618.9$ ft), and the smallest flat cut that secures 618 ft is $1-\rho^{*}=7.2\%$ of $\Pbar$ by Theorem~\ref{thm:margin}---a pumping level of $261.8$, which is $31.5\%$ below the 1991--2023 mean of $382.1$. The Stage-I rule settles at the flat-80\% attractor ($622.0$ ft) because the head it converges to lies below 660 ft where its cut is active, and the cascade settles at the flat-60\% attractor ($628.4$ ft) in its 40\% stage. Under the 5th- and 10th-percentile floors every rule including BAU clears 618 ft.

\begin{table}[t]
\centering
\caption{Edwards Aquifer: attractors, kernels and open-loop supply by policy.}
\label{tab:edw}
\small
\resizebox{\textwidth}{!}{%
\begin{tabular}{lrrrllrr}
\toprule
Policy & \multicolumn{3}{c}{attractor $H^*$ (ft)} & 618-ft kernel & 660-ft $b_1/b_2/b_3$ & supply & cut active \\ & worst & q05 & q10 & (worst floor) & (worst floor) & &  \\
\midrule
BAU ($\bar P$) & 615.7 & 625.3 & 626.3 & empty from $T=13$ & 675.1 / 695.3 / $\emptyset$ & 282.2 & 0\% \\
flat 90\% & 618.9 & 628.5 & 629.5 & whole safe set & 674.0 / 692.8 / $\emptyset$ & 253.9 & 100\% \\
flat 80\% & 622.0 & 631.6 & 632.6 & whole safe set & 672.9 / 690.2 / $\emptyset$ & 225.7 & 100\% \\
flat 70\% & 625.2 & 634.8 & 635.8 & whole safe set & 671.8 / 687.7 / 709.0 & 197.5 & 100\% \\
flat 60\% & 628.4 & 638.0 & 638.9 & whole safe set & 670.8 / 685.2 / 704.5 & 169.3 & 100\% \\
flat 50\% & 631.5 & 641.1 & 642.1 & whole safe set & 669.7 / 682.7 / 700.1 & 141.1 & 100\% \\
zero pumping & 647.3 & 656.9 & 657.9 & whole safe set & 664.3 / 670.1 / 677.8 & 0.0 & 100\% \\
Stage I ($-20\%$ below 660 ft) & 622.0 & 631.6 & 632.6 & whole safe set & 675.1 / 695.3 / $\emptyset$ & 262.4 & 35\% \\
CPM cascade & 628.4 & 636.4 & 637.4 & whole safe set & 675.1 / 695.3 / $\emptyset$ & 254.9 & 35\% \\
\bottomrule
\end{tabular}
}
\par\smallskip\begin{minipage}{0.97\linewidth}\footnotesize Attractor $H^*=(\alpha+\beta F+\gamma P)/(1-a)$ of the stage that contains it (self-consistent fixed point of the closed loop) under the three recharge floors $F$. 618-ft kernel under the worst floor: whole safe set iff $H^*\ge 618$ (Corollary~\ref{cor:edw}); otherwise empty beyond $T_{\mathrm{empty}}$. 660-ft column: lower boundary at $T=1/2/3$ under the worst floor; zero pumping is the only rule whose 660-ft kernel survives to $T=3$, and none survives $T=\infty$ because $H^*_0=(\alpha+\beta F)/(1-a)=647.3$ ft $<660$ (Theorem~\ref{thm:margin}). Supply: mean prescribed pumping on the observed 1934--1990 heads ($10^3$ acre-ft yr$^{-1}$) and share of years with the cut active.\end{minipage}
\end{table}


The 618-ft kernels follow. For every cut policy they are the whole safe set at every horizon under every floor. For BAU under the worst floor they are the empty-in-finite-time case of Theorem~\ref{thm:kernel}(d): $b_T=\Hstar+(618-\Hstar)/a^{T}$ runs $618.8$, $619.8$, $621.2$, $625.6$, $639.5$, $658.4$ and $692.4$ ft at $T=1,2,3,5,8,10,12$ and passes the 710-ft ceiling at $T=13$, in agreement with the closed form $\Tempty=\lfloor12.70\rfloor+1$. A manager who pumps the historical mean through a repeat of 1956 recharge can keep J-17 above 618 ft for twelve years only by starting above $692$ ft, which is above the highest annual mean in the ninety-year record ($692.0$ ft). Finite-duration floors (SI Section~S10) say the same with less drama: five consecutive 1956 years under BAU are absorbed from $625.6$ ft, ten from $658.4$ ft, fifteen from nowhere; every cut policy absorbs any run.

The 660-ft trigger is a different object. Theorem~\ref{thm:margin} asks whether the zero-pumping attractor clears it: $\Hstar_0=(\alpha+\beta\underline R)/(1-a)=647.3$, $656.9$ and $657.9$ ft under the three floors, short of 660 by $12.7$, $3.1$ and $2.1$ ft. No admissible rule can hold 660 ft from itself under any declared floor; the kernels are empty in finite time for every rule, at $T=3$ for BAU, the Stage-I rule and the cascade (whose 660-ft kernels coincide, since all three pump $\Pbar$ at and above 660 ft), at $T=6$ for zero pumping (SI Section~S8). The institutional trigger is a line the aquifer is managed \emph{around}, not a line that can be certified; the physical trigger is one that can.

\begin{figure}[t]\centering
\includegraphics[width=\textwidth]{fig_edw_kernels.pdf}
\caption{Edwards Aquifer kernels. (a)~Worst-floor attractors $\Hstar$ of the rule family under the three recharge floors (Corollary~\ref{cor:edw}, Proposition~\ref{prop:stage}), with 90\% intervals from the archived residual bootstrap of the fit; the staged rules inherit the interval of the flat cap of their stage at the attractor. (b)~Lower boundaries of the 618-ft kernels under the 1956 floor: Lipschitz-certified kernels with the declared defect for five rules (solid), the nominal BAU kernel (dashed), and the exact bounded-defect BAU kernel under the out-of-sample defect of $21.8\ft$ (dotted). The dotted grey line is the 710-ft ceiling.}
\label{fig:edw-kernels}
\end{figure}


\subsection{Supply and retention}\label{sec:edw-supply}

Protection ties are broken by supply. On the observed 1934--1990 heads the Stage-I rule's cut is active in 20 of 57 years ($35\%$) and it prescribes a mean of $262.36$, against $253.94$ for the flat 90\% cap and $225.73$ for flat 80\%; the observed heads lie below 660, 650, 640 and 630 ft in 20, 12, 5 and 1 of 57 years and the cascade prescribes $254.93$, against $169.29$ for the flat 60\% cap whose protection it matches. Clause (a) of the retention rule holds for both staged rules at every reading (they are never less protective than BAU), and clause (b) holds at the 618-ft worst-floor reading: the Stage-I rule supplies $3.3\%$ more than the least protective flat cap that matches it at the nominal level (flat 90\%) and $16.2\%$ more than the one that matches it at the certified level (flat 80\%); the cascade supplies $50.6\%$ more than flat 60\% and $0.4\%$ more than flat 90\%. Both staged rules are retained. On the 33 out-of-sample heads the cut is active in 11 years ($33\%$) and the supplies are $263.35$ and $258.22$ ($10^3$ acre-ft yr$^{-1}$).

Open-loop supply evaluates a rule on heads it did not produce. The closed-loop replay (model driven by observed recharge, pumping by the rule, from the observed 1934 head; Figure~\ref{fig:edw-replay}a) corrects for this: the Stage-I rule supplies $258.98$ with the cut active in $41\%$ of years, $2.0\%$ more than flat 90\% and $14.7\%$ more than flat 80\%; the cascade supplies $254.70$, $50.4\%$ more than flat 60\%. The retention conclusions survive the correction. Out of sample (from the 1990 head through the 33 transitions to 2023; Figure~\ref{fig:edw-replay}b) the Stage-I rule supplies $268.48$ with the cut active in $24\%$ of years and ends 2023 at $643.3$ ft against an observed $635.7$ ft, the observed series having been pumped at a mean of $382$ rather than $282$. Under a recharge pinned at the 1956 floor for 57 years the staged rules spend most of the time in their deepest stage and their supply advantage shrinks to the stage margin ($226.7$ against $225.7$ for Stage I and flat 80\%; $174.5$ against $169.3$ for the cascade and flat 60\%; SI Section~S10).

\subsection{Certified horizons and the bounded-defect audit}\label{sec:edw-horizon}

With $a<1$ the Lipschitz accumulation saturates: $r_T=\varepsilon(1-a^{T})/(1-a)$ runs $15.41$, $26.91$, $35.49$, $41.89$, $46.66$ ft at $T=1,\dots,5$ and tends to $r_\infty=60.70$ ft. The certified 618-ft kernel at horizon $T$ is the nominal kernel of $618+r_T$, and by Corollary~\ref{cor:edw} its boundary is $\Hstar+(618+r_T-\Hstar)/a^{T}$ whenever that exceeds $618+r_T$. For BAU under the worst floor the certified boundaries are $639.4$, $668.2$ and $706.7$ ft at $T=1,2,3$ and the kernel is empty at $T=4$; the Stage-I rule and the cascade reach the same $706.7$ ft at $T=3$ (their attractors differ, but the $1/a^{T}$ factor dominates), flat 80\% reaches $697.8$ ft, and only zero pumping survives to $T=4$ ($687.9$ ft) under the worst floor; under the two quantile floors the 50\% and 60\% caps reach $T=4$ and zero pumping $T=5$. The registered Lipschitz certificate for 660 ft fails at $T=2$ for every rule. Table~\ref{tab:horizons} collects the horizons and Figure~\ref{fig:edw-kernels} draws the boundaries.

Proposition~\ref{prop:lip}(ii) says that for an affine closed-loop map the exact bounded-defect kernel of Theorem~\ref{thm:exact} coincides with the Lipschitz certificate wherever the backward recursion binds. The BAU boundaries are $639.43$, $668.16$, and $706.66$ ft, identical to the certificate, with the kernel empty at $T=4$; the positive flat caps likewise coincide at every reported horizon, as does zero pumping from $T=3$ onward. The departures are of two kinds and neither lengthens a horizon. For zero pumping at $T=1,2$ the certificate's inflation of the initial state is the binding constraint and the exact kernel is larger ($628.7$ and $643.0$ against $633.4$ and $644.9$ ft). For the two staged rules the closed-loop map is piecewise affine rather than affine, because a perturbed head can cross 660 ft and engage the cut, and the exact kernel is larger by $2$--$7$ ft at $T=2,3$ (Stage I $702.8$ against $706.7$ ft at $T=3$; cascade $663.1$ and $699.9$ against $665.3$ and $706.7$ ft) while still emptying at $T=4$. These departures sharpen boundary locations but add no certified year: for positive pumping the three-year horizon is a property of the declared class, not a looseness in the Lipschitz method. Proposition~\ref{prop:finite}(ii) makes the point structurally: an indefinite certificate requires an attractor at least $r_\infty=60.7$ ft above the threshold, i.e.\ $678.7$ ft, and the highest attractor in the family is $657.9$ ft (zero pumping, q10 floor).

Because the out-of-sample residual exceeds the declaration, the audit of Table~\ref{tab:edw-unc} asks how large a uniform defect the 618-ft kernel can carry at each horizon. The supremum at $T=3$ is $16.01$ ft for BAU, $16.71$ ft for Stage I, $17.23$ ft for the cascade, $17.62$ ft for flat 80\% and $24.04$ ft for zero pumping; at $T=4$ it is $9.91$--$11.60$ ft for the pumping rules and $17.93$ ft for zero pumping. The declared $15.41$ ft clears every $T=3$ supremum and no $T=4$ supremum except zero pumping's; the out-of-sample $21.81$ ft clears only zero pumping's at $T=3$: every positive-pumping rule's exact kernel is empty at $T=3$ (BAU boundaries $648.0$ and $688.2$ ft at $T=1,2$), while zero pumping's kernel is empty at $T=4$. The out-of-sample defect accumulates to $r_\infty=85.9$ ft. For positive-pumping rules the certified horizon is three years on the declared defect and two on the audited one; zero pumping adds one year in each case (four and three years, respectively).

\begin{table}[t]
\centering
\caption{Edwards Aquifer: supply, protection and the bounded-defect audit.}
\label{tab:edw-unc}
\small
\setlength{\tabcolsep}{4pt}
\begin{tabular}{lrrrr}
\toprule
\multicolumn{5}{l}{\emph{Mean annual supply $P$ ($10^3$ acre-ft yr$^{-1}$) and $P(H^*\ge618)$}} \\
\midrule
policy & \shortstack{open loop\\1934--90} & \shortstack{closed loop\\fit 1934--90} & \shortstack{closed loop OOS\\1991--2023} & \shortstack{$P(H^*\ge618)$\\archived / rerun} \\
\midrule
BAU ($\bar P$) & 282.16 & 282.16 (0\%) & 282.16 (0\%) & 0.41 / 0.40 \\
flat 90\% & 253.94 & 253.94 (100\%) & 253.94 (100\%) & 0.53 / 0.52 \\
flat 80\% & 225.73 & 225.73 (100\%) & 225.73 (100\%) & 0.65 / 0.64 \\
flat 60\% & 169.29 & 169.29 (100\%) & 169.29 (100\%) & 0.83 / 0.82 \\
zero pumping & 0.00 & 0.00 (100\%) & 0.00 (100\%) & 0.99 / 0.99 \\
Stage I ($-20\%$ below 660 ft) & 262.36 & 258.98 (41\%) & 268.48 (24\%) & 0.65 / -- \\
CPM cascade & 254.93 & 254.70 (41\%) & 265.06 (24\%) & 0.83 / -- \\
\midrule
\multicolumn{5}{l}{\emph{Exact-kernel supremum defect $\bar\varepsilon_T$ (ft) by horizon}} \\
\midrule
policy & $T=2$ & $T=3$ & $T=4$ & $T=5$ \\
\midrule
BAU ($\bar P$) & 28.75 & 16.01 & 9.91 & 6.44 \\
flat 90\% & 29.55 & 16.82 & 10.71 & 7.25 \\
flat 80\% & 30.36 & 17.62 & 11.51 & 8.05 \\
Stage I ($-20\%$ below 660 ft) & 29.67 & 16.71 & 10.94 & 7.67 \\
CPM cascade & 29.67 & 17.23 & 11.60 & 8.49 \\
zero pumping & 36.77 & 24.04 & 17.93 & 14.47 \\
\bottomrule
\end{tabular}
\par\smallskip\begin{minipage}{0.97\linewidth}\footnotesize Closed-loop replays use observed recharge and rule-based pumping, starting from 1934 (fit) or 1990 (out of sample); percentages in parentheses are years with the cut active. $P(H^*\ge618)$ is the share of 5,000 bootstrap refits whose worst-floor attractor clears 618 ft (archived / independent seeds 20260905 / 20261015). At the declared 15.41-ft defect, every rule certifies three years and only zero pumping lasts four; at the out-of-sample 21.81-ft defect, only zero pumping certifies three years.\end{minipage}
\end{table}


\subsection{Replays and the stochastic reading}\label{sec:edw-stoch}

The 1950s drought is the only episode in the record that reached the physical trigger's neighbourhood: J-17 averaged $623.15$ ft in 1956 (daily minimum $612.51$ ft). Replaying 1951--1956 from the observed 1950 head with actual recharge and each rule's pumping, the model reaches $629.7$ ft under BAU, $632.3$ under flat 90\%, $634.9$ under Stage I and $637.1$ ft under the cascade, all above 618; with the actual (larger) pumping of the 1950s the model reaches $631.2$ ft against the observed $623.2$, a bias of $8.1$ ft that is within the declared defect and in the unsafe direction. Under the fit-window record the staged rules keep the model above $632.9$ and $636.2$ ft; under 1991--2023 recharge all rules end 2023 between $639$ and $671$ ft (SI Section~S10).

The stochastic reading resamples fit residuals, refits \eqref{eq:edw-map} 5000 times and asks how often the worst-floor attractor clears 618 ft: $0.41$ for BAU, $0.53$ for flat 90\%, $0.65$ for flat 80\% and Stage I, $0.83$ for flat 60\% and the cascade, $0.99$ for zero pumping, with a bootstrap standard deviation of $10$--$13$ ft on every attractor; an independent re-run with a different seed reproduces each probability to within $0.012$ (SI Section~S11). The flat-90\% minus BAU attractor gap, $3.2$ ft at the point estimate, has a 90\% bootstrap interval of $[1.9,5.1]$ ft and is positive in every replicate (minimum $0.5$ ft), because $\gamma<0$ in every refit. The ordering of the rules is sampling-stable; the certification of 618 ft by a 7.2\% cut is a statement about the registered fit, and at the fit's own sampling uncertainty it is a coin toss for BAU and a two-to-one bet for the Stage-I rule.

\begin{figure}[t]\centering
\includegraphics[width=\textwidth]{fig_edw_replay.pdf}
\caption{Closed-loop replays of the Edwards map (observed recharge, pumping prescribed by the rule from the model head, 710-ft ceiling). (a)~From the observed 1934 head through 1990; (b)~from the observed 1990 head through 2023. The observed series was pumped at a mean of $282$ ($10^3$ acre-ft) before 1991 and $382$ after; the replays pump $\Pbar=282$ or the rule's fraction of it.}
\label{fig:edw-replay}
\end{figure}
